% =========================================================
% 05_matrices_transformations.tex
% Vectors and Linear Algebra
% =========================================================

% =========================================================
% SECTION DIVIDER
% =========================================================

\section{Matrices and Linear Transformations}

% =========================================================
% MATRICES
% =========================================================

\begin{frame}{Matrices}

A matrix is a rectangular array of numbers.

\[
A=
\begin{bmatrix}
a_{11} & a_{12}\\
a_{21} & a_{22}
\end{bmatrix}.
\]

\vspace{0.4cm}

A matrix with $m$ rows and $n$ columns has dimensions

\[
\boxed{m\times n}.
\]

\vspace{0.4cm}

Matrices can represent:

\begin{itemize}
    \item systems of linear equations
    \item transformations
    \item datasets
    \item relationships between variables
\end{itemize}

\end{frame}


% =========================================================
% MATRIX-VECTOR MULTIPLICATION
% =========================================================

\begin{frame}{Matrix--Vector Multiplication}

Let

\[
A=
\begin{bmatrix}
a_{11} & a_{12}\\
a_{21} & a_{22}
\end{bmatrix},
\qquad
\vect{x}
=
\begin{bmatrix}
x_1\\
x_2
\end{bmatrix}.
\]

Then

\[
A\vect{x}
=
\begin{bmatrix}
a_{11}x_1+a_{12}x_2\\
a_{21}x_1+a_{22}x_2
\end{bmatrix}.
\]

\vspace{0.4cm}

Equivalently,

\[
\boxed{
A\vect{x}
=
x_1
\begin{bmatrix}
a_{11}\\
a_{21}
\end{bmatrix}
+
x_2
\begin{bmatrix}
a_{12}\\
a_{22}
\end{bmatrix}
}
\]

So matrix--vector multiplication is a linear combination
of the columns of $A$.

\end{frame}


% =========================================================
% MATRIX AS A TRANSFORMATION
% =========================================================

\begin{frame}{A Matrix as a Transformation}

A matrix can define a function
\vspace{-0.1cm}
\[
T:\mathbb{R}^n\rightarrow\mathbb{R}^m
\]

through
\[
\boxed{
T(\vect{x})=A\vect{x}.
}
\]

The matrix $A$ determines how the input vector is transformed.

\[
\boxed{
\vect{x}
\quad\xrightarrow{\quad A\quad}\quad
A\vect{x}
}
\]

If

\[
A\in\mathbb{R}^{m\times n},
\]

then

\[
\vect{x}\in\mathbb{R}^n
\quad\Longrightarrow\quad
A\vect{x}\in\mathbb{R}^m.
\]

\end{frame}


% =========================================================
% LINEAR TRANSFORMATION
% =========================================================

\begin{frame}{Linear Transformations}

A transformation $T$ is linear if, for all vectors $\vect{u},\vect{v}$
and scalars $c$,

\[
\boxed{
T(\vect{u}+\vect{v})
=
T(\vect{u})+T(\vect{v})
}
\]

and

\[
\boxed{
T(c\vect{u})
=
cT(\vect{u}).
}
\]

These two properties can be combined as

\[
\boxed{
T(c_1\vect{u}+c_2\vect{v})
=
c_1T(\vect{u})+c_2T(\vect{v}).
}
\]

\end{frame}


% =========================================================
% EXAMPLE OF TRANSFORMATION
% =========================================================

\begin{frame}{Example: Scaling Transformation}

Consider
\vspace{-0.1cm}
\[
A=
\begin{bmatrix}
2&0\\
0&3
\end{bmatrix}.
\]

For
\vspace{-0.1cm}
\[
\vect{x}
=
\begin{bmatrix}
x\\
y
\end{bmatrix},
\]

we obtain

\[
A\vect{x}
=
\begin{bmatrix}
2x\\
3y
\end{bmatrix}.
\]

Therefore, the transformation

\[
T(\vect{x})=A\vect{x}
\]

scales the $x$-coordinate by $2$ and the $y$-coordinate
by $3$.

\end{frame}


% =========================================================
% IDENTITY TRANSFORMATION
% =========================================================

\begin{frame}{Identity Transformation}

The identity matrix is

\[
I=
\begin{bmatrix}
1&0\\
0&1
\end{bmatrix}.
\]

For any vector $\vect{x}$,

\[
\boxed{
I\vect{x}=\vect{x}.
}
\]

The corresponding transformation

\[
T(\vect{x})=I\vect{x}
\]

leaves every vector unchanged.
\vspace{0.1cm}

Thus, the identity transformation acts as the
\textbf{neutral transformation} under composition.

\end{frame}


% =========================================================
% ROTATION
% =========================================================

\begin{frame}{Rotation as a Linear Transformation}

A rotation through angle $\theta$ in $\mathbb{R}^2$ is represented by

\[
\boxed{
R_\theta
=
\begin{bmatrix}
\cos\theta & -\sin\theta\\
\sin\theta & \cos\theta
\end{bmatrix}.
}
\]

Therefore,

\[
T(\vect{x})
=
R_\theta\vect{x}.
\]

\vspace{0.4cm}

For $\theta=90^\circ$,

\[
R_{90^\circ}
=
\begin{bmatrix}
0&-1\\
1&0
\end{bmatrix}.
\]

\vspace{0.3cm}

The transformation rotates every vector by $90^\circ$.

\end{frame}


% =========================================================
% REFLECTION
% =========================================================

\begin{frame}{Reflection as a Linear Transformation}

Reflection about the $x$-axis can be represented by

\[
\boxed{
A=
\begin{bmatrix}
1&0\\
0&-1
\end{bmatrix}.
}
\]

For

\[
\vect{x}
=
\begin{bmatrix}
x\\
y
\end{bmatrix},
\]

we obtain

\[
A\vect{x}
=
\begin{bmatrix}
x\\
-y
\end{bmatrix}.
\]

\vspace{0.4cm}

The $x$-coordinate remains unchanged while the
$y$-coordinate changes sign.

\end{frame}


% =========================================================
% MATRIX COLUMNS AS TRANSFORMED BASIS VECTORS
% =========================================================

\begin{frame}{Columns Determine the Transformation}

Let
\vspace{-0.1cm}
\[
A=
\begin{bmatrix}
|&|\\
\vect{a}_1&\vect{a}_2\\
|&|
\end{bmatrix}.
\]

For
\vspace{-0.1cm}
\[
\vect{x}
=
\begin{bmatrix}
x_1\\
x_2
\end{bmatrix},
\]

we have
\vspace{-0.1cm}
\[
A\vect{x}
=
x_1\vect{a}_1+x_2\vect{a}_2.
\]

In particular,
\vspace{-0.1cm}
\[
A
\begin{bmatrix}
1\\
0
\end{bmatrix}
=
\vect{a}_1,
\qquad
A
\begin{bmatrix}
0\\
1
\end{bmatrix}
=
\vect{a}_2.
\]

Thus, the columns of a matrix are the images of the
standard basis vectors.

\end{frame}


% =========================================================
% COMPOSITION
% =========================================================

\begin{frame}{Composition of Transformations}

Suppose two linear transformations are represented by

\[
T(\vect{x})=A\vect{x}
\]

and

\[
S(\vect{x})=B\vect{x}.
\]

Applying $T$ first and then $S$ gives
\vspace{-0.1cm}
\[
S(T(\vect{x}))
=
B(A\vect{x}).
\]

Therefore,
\vspace{-0.1cm}
\[
\boxed{
S\circ T
\quad\longleftrightarrow\quad
BA.
}
\]
\vspace{-0.1cm}

Matrix multiplication therefore provides a compact way to
represent the \textbf{composition of linear transformations}.

\end{frame}


% =========================================================
% ORDER MATTERS
% =========================================================

\begin{frame}{Order of Composition Matters}

In general,

\[
\boxed{
AB\neq BA.
}
\]

Therefore,

\[
B(A\vect{x})
\neq
A(B\vect{x})
\]

in general.

\vspace{0.5cm}

This means that applying transformation $A$ followed by $B$
is generally different from applying $B$ followed by $A$.

\vspace{0.4cm}

\begin{block}{Key Idea}
Matrix multiplication is generally \textbf{not commutative}.
\end{block}

\end{frame}


% =========================================================
% COMPOSITION EXAMPLE
% =========================================================

\begin{frame}{Composition: Example}

Consider

\[
A=
\begin{bmatrix}
2&0\\
0&1
\end{bmatrix},
\qquad
B=
\begin{bmatrix}
1&0\\
0&3
\end{bmatrix}.
\]

Then

\[
BA
=
\begin{bmatrix}
2&0\\
0&3
\end{bmatrix}.
\]

For this particular pair,

\[
AB=BA.
\]

But this is not generally true.

For transformations such as rotations, reflections,
and shears, the order can produce different results.

\end{frame}


% =========================================================
% THREE-DIMENSIONAL TRANSFORMATIONS
% =========================================================

\begin{frame}{Three-Dimensional Linear Transformations}

A linear transformation in three dimensions has the form

\[
\boxed{
T:\mathbb{R}^3\rightarrow\mathbb{R}^3
}
\]

with
\vspace{-0.1cm}
\[
T(\vect{x})=A\vect{x},
\]

where
\vspace{-0.1cm}
\[
A=
\begin{bmatrix}
a_{11}&a_{12}&a_{13}\\
a_{21}&a_{22}&a_{23}\\
a_{31}&a_{32}&a_{33}
\end{bmatrix}.
\]

The matrix can represent transformations such as:

\begin{itemize}
    \item scaling
    \item rotation
    \item reflection
    \item shear
\end{itemize}

\end{frame}


% =========================================================
% 3D SCALING
% =========================================================

\begin{frame}{3D Scaling}

A three-dimensional scaling transformation can be written as

\[
A=
\begin{bmatrix}
s_x&0&0\\
0&s_y&0\\
0&0&s_z
\end{bmatrix}.
\]

Then

\[
\begin{bmatrix}
x'\\
y'\\
z'
\end{bmatrix}
=
A
\begin{bmatrix}
x\\
y\\
z
\end{bmatrix}
\]

gives
\vspace{-0.1cm}
\[
\boxed{
x'=s_xx,\qquad
y'=s_yy,\qquad
z'=s_zz.
}
\]

Different scaling factors can independently stretch or
compress the three coordinate directions.

\end{frame}


% =========================================================
% 3D ROTATION
% =========================================================

\begin{frame}{Rotation About the $z$-Axis}

A rotation through $\theta$ about the $z$-axis is represented by

\[
\boxed{
R_z(\theta)
=
\begin{bmatrix}
\cos\theta&-\sin\theta&0\\
\sin\theta&\cos\theta&0\\
0&0&1
\end{bmatrix}.
}
\]

\vspace{0.4cm}

The $x$ and $y$ coordinates rotate in their plane while
the $z$ coordinate remains unchanged.

\vspace{0.4cm}

This type of transformation is fundamental in

\[
\text{3D graphics}
\qquad\text{and}\qquad
\text{computer vision}.
\]

\end{frame}


% =========================================================
% TRANSFORMATION PIPELINE
% =========================================================

\begin{frame}{Transformation Pipeline}

Multiple transformations can be combined into a single matrix.

\[
\vect{x}
\xrightarrow{\ A\ }
A\vect{x}
\xrightarrow{\ B\ }
BA\vect{x}
\xrightarrow{\ C\ }
CBA\vect{x}.
\]

Therefore,

\[
\boxed{
T(\vect{x})=CBA\vect{x}.
}
\]

\vspace{0.5cm}

This idea is central to transformation pipelines in:

\begin{itemize}
    \item computer graphics
    \item robotics
    \item computer vision
    \item geometric modelling
\end{itemize}

\end{frame}


% =========================================================
% HOMOGENEOUS COORDINATES PREVIEW
% =========================================================

\begin{frame}{Beyond Linear Transformations}

Pure linear transformations preserve the origin:

\[
\boxed{
T(\mathbf{0})=\mathbf{0}.
}
\]

Translations do not satisfy this property.

\vspace{0.4cm}

In computer graphics, translations are commonly incorporated
using \textbf{homogeneous coordinates}.

For example,

\[
\begin{bmatrix}
x'\\
y'\\
1
\end{bmatrix}
=
\begin{bmatrix}
1&0&t_x\\
0&1&t_y\\
0&0&1
\end{bmatrix}
\begin{bmatrix}
x\\
y\\
1
\end{bmatrix}.
\]

This allows translation and linear transformations to be
combined into matrix operations.

\end{frame}